\documentclass[11pt,a4paper]{article}

\usepackage[a4paper,margin=25mm]{geometry}
\usepackage{fontspec}
\usepackage{polyglossia}
\setmainlanguage{english}
\setmainfont{DejaVu Serif}
\setsansfont{DejaVu Sans}
\setmonofont{DejaVu Sans Mono}

\usepackage{amsmath,amssymb,amsfonts,mathtools}
\usepackage{booktabs,longtable,array,tabularx,multirow}
\usepackage{enumitem}
\usepackage{hyperref}
\usepackage{xcolor}
\usepackage{listings}
\usepackage{tikz}
\usetikzlibrary{positioning,arrows.meta,shapes.geometric}
\usepackage{caption}
\usepackage{float}
\usepackage{fancyhdr}
\usepackage{titlesec}
\usepackage{tocloft}
\usepackage{parskip}

\hypersetup{
    colorlinks=true,
    linkcolor=blue!60!black,
    urlcolor=blue!60!black,
    citecolor=blue!60!black,
    pdftitle={dnipro Reference Manual},
    pdfauthor={dnipro project}
}

\setlength{\headheight}{14pt}
\pagestyle{fancy}
\fancyhf{}
\lhead{\texttt{dnipro} Reference Manual}
\rhead{v0.1.0}
\cfoot{\thepage}

\lstdefinelanguage{spice}{
    morekeywords={DC,SIN,PULSE,PWL,SQUARE,SAW},
    morecomment=[l]{*},
    morestring=[b]",
    sensitive=false
}

\lstdefinestyle{code}{
    basicstyle=\ttfamily\small,
    breaklines=true,
    frame=single,
    columns=fullflexible,
    keepspaces=true,
    showstringspaces=false,
    backgroundcolor=\color{gray!5},
    rulecolor=\color{gray!40}
}

\lstdefinestyle{shell}{
    basicstyle=\ttfamily\small,
    breaklines=true,
    frame=single,
    columns=fullflexible,
    keepspaces=true,
    showstringspaces=false,
    backgroundcolor=\color{black!4},
    rulecolor=\color{gray!60}
}

\lstset{style=code}

\titleformat{\section}{\Large\bfseries}{\thesection}{1em}{}
\titleformat{\subsection}{\large\bfseries}{\thesubsection}{1em}{}
\titleformat{\subsubsection}{\normalsize\bfseries}{\thesubsubsection}{1em}{}

\newcommand{\project}{\texttt{dnipro}}
\newcommand{\code}[1]{\texttt{#1}}
\newcommand{\mna}{Modified Nodal Analysis}
\newcommand{\sympy}{\texttt{SymPy}}
\newcommand{\ngspice}{\texttt{ngspice}}
\newcommand{\pyver}{Python~3.11+}

% Signature box for API entries
\newenvironment{apisig}{%
  \begin{trivlist}\item\relax\ttfamily\small
}{%
  \end{trivlist}
}

\begin{document}

% ============================================================
% Title page
% ============================================================
\begin{titlepage}
    \centering
    \vspace*{20mm}
    {\Huge\bfseries dnipro\par}
    \vspace{8mm}
    {\Large Symbolic Equation-Based MNA Framework\par}
    \vspace{3mm}
    {\large for SPICE-like Circuit Modeling in Python and SymPy\par}
    \vspace{18mm}
    {\Large\bfseries Reference Manual\par}
    \vspace{5mm}
    {\large Version 0.1.0\par}
    \vfill
    \begin{tabular}{ll}
        Package name:            & \texttt{dnipro} \\
        Implementation language: & Python 3.11+ \\
        Core symbolic engine:    & SymPy \(\geq\) 1.14 \\
        External verifier:       & ngspice (optional) \\
        Entry point:             & \texttt{dnipro analyze} \\
    \end{tabular}
    \vfill
    {\large \today\par}
\end{titlepage}

\tableofcontents
\newpage

% ============================================================
% PART I — USER GUIDE
% ============================================================
\part{User Guide}

% ------------------------------------------------------------
\section{Introduction}
\label{sec:introduction}

\project{} is a \pyver{} library and command-line tool for \emph{symbolic circuit analysis}
based on equation-form \mna{} (MNA).
Given a SPICE-like netlist, \project{} constructs an algebraic system of residual equations
in \sympy{} symbolic form:
\begin{equation}
    F\!\left(x(t),\,\dot{x}(t),\,t,\,p\right) = 0,
\end{equation}
where \(x(t)\) collects node voltages and branch currents,
\(\dot{x}(t)\) denotes their time derivatives,
\(t\) is time, and \(p\) is the set of symbolic or numeric circuit parameters.

\subsection{What dnipro is}

\project{} is designed for use cases where the \emph{equations themselves} are the output:
teaching circuit theory, generating symbolic transfer functions, preparing
derivations for technical reports, or cross-checking hand analysis against a reference tool.
The library exports residual equations as human-readable text, JSON (for downstream
tooling), or LaTeX (for direct inclusion in documents). For linear systems it can
also derive and report the classical algebraic MNA form \(A x = b\) and the
time-domain descriptor form \(Gx + C\dot{x} = b\).

\subsection{What dnipro is not}

\project{} is not a SPICE replacement.
It does not perform transient numerical integration, does not compute operating
points automatically for nonlinear circuits, and does not claim full SPICE netlist
compatibility.
External \ngspice{} is used only as a verification oracle in integration tests
and tutorial comparison artifacts. It is never required by the runtime library.

\subsection{Relationship to the specification}

The companion document \emph{SPEC\_dnipro\_v1\_symbolic\_mna\_framework} defines the
version~1 requirements.
This manual describes the implemented library and CLI as they exist in \project{}~0.1.0.

\subsection{Relationship to the User's Guide}

This reference manual is organized around the implemented API, syntax, CLI, and
export formats. The companion document \emph{Dnipro User's Guide by Examples}
is organized around runnable circuits in \code{examples/user\_guide}. Use the
User's Guide when learning the system from concrete circuits, \project{} Python
scripts, and matching \ngspice{} templates.

% ------------------------------------------------------------
\section{Installation and Setup}
\label{sec:installation}

\subsection{Requirements}

\begin{itemize}[noitemsep]
    \item Python 3.11 or later
    \item \sympy{} 1.14 or later (installed automatically as a dependency)
    \item \ngspice{} — optional; required only for integration tests marked
          \code{@pytest.mark.ngspice}
    \item NumPy --- optional; used by selected tutorial scripts for fast numeric
          AC sweeps and comparison with \ngspice{} output.
\end{itemize}

\subsection{Installing from source}

\begin{lstlisting}[style=shell]
git clone <repository>
cd dnipro
pip install -e .
\end{lstlisting}

After installation the command \code{dnipro} is available on the \code{PATH}
and the package \code{dnipro} is importable from Python.

\subsection{Verifying the installation}

\begin{lstlisting}[style=shell]
dnipro --help
python -c "import dnipro; print(dnipro.__version__)"
\end{lstlisting}

\subsection{Running the test suite}

\begin{lstlisting}[style=shell]
pip install pytest
pytest
\end{lstlisting}

To include \ngspice{}-dependent integration tests:

\begin{lstlisting}[style=shell]
pytest -m ngspice
\end{lstlisting}

Tests marked \code{ngspice} are automatically skipped when the \ngspice{} binary
is absent from \code{PATH}.

% ------------------------------------------------------------
\section{Quick Start}
\label{sec:quickstart}

\subsection{Example 1: Resistive voltage divider (time domain)}

Create a netlist file \code{divider.cir}:

\begin{lstlisting}[language=spice]
* Resistive voltage divider
V1 in 0 Vin
R1 in out R1
R2 out 0 R2
\end{lstlisting}

Analyze it:

\begin{lstlisting}[style=shell]
dnipro analyze divider.cir
\end{lstlisting}

Sample output:

\begin{lstlisting}
dnipro residual system
domain: time
nodes: in, out
branch currents: V1
unknowns:
  [0] v_in(t)
  [1] v_out(t)
  [2] i_V1(t)
residuals:
  KCL(in): (v_in(t) - v_out(t))/R1 + i_V1(t) = 0
  KCL(out): -(v_in(t) - v_out(t))/R1 + (v_out(t))/R2 = 0
  V1: v_in(t) - Vin = 0
\end{lstlisting}

\subsection{Example 2: RC low-pass filter (Laplace domain)}

\begin{lstlisting}[language=spice]
* RC low-pass filter
V1 in 0 Vin
R1 in out R
C1 out 0 C
\end{lstlisting}

\begin{lstlisting}[style=shell]
dnipro analyze rc.cir --domain laplace --output latex
\end{lstlisting}

The output is a \LaTeX{} fragment containing the Laplace-domain residuals,
suitable for inclusion in a report with \verb|\input{...}|.

\subsection{Example 3: Numeric substitution}

\begin{lstlisting}[style=shell]
dnipro analyze rc.cir --domain laplace --numeric \
    --param Vin=1 --param R=1k --param C=1u \
    --output json
\end{lstlisting}

All three symbolic parameters (\code{Vin}, \code{R}, \code{C}) are substituted
before export.
If any required parameter is missing \project{} prints a clear error and exits with
status~1.

\subsection{Example 4: Classical linear MNA report}

For a linear circuit, request the derived algebraic matrix form:

\begin{lstlisting}[style=shell]
dnipro analyze divider.cir --linear-system
\end{lstlisting}

Sample output includes row names, the unknown vector, the coefficient matrix,
and the right-hand-side vector:

\begin{lstlisting}
dnipro linear MNA system
domain: time
form: A*x = b

row names:
  [0] KCL(in)
  [1] KCL(out)
  [2] V1

unknown vector x:
  [0] v_in(t)
  [1] v_out(t)
  [2] i_V1(t)

A matrix:
  [1/R1, -1/R1, 1]
  [-1/R1, 1/R1 + 1/R2, 0]
  [1, 0, 0]

b vector:
  [0]
  [0]
  [Vin]
\end{lstlisting}

\subsection{Example 5: Descriptor MNA report}

For a time-domain linear circuit with reactive elements, request descriptor
form:

\begin{lstlisting}[style=shell]
dnipro analyze rc.cir --descriptor-system --output latex
\end{lstlisting}

The output is a \LaTeX{} fragment containing the row mapping, unknown and
derivative mapping, and displayed \(G\), \(C\), \(x\), \(\dot{x}\), and \(b\)
objects for the equation \(Gx + C\dot{x} = b\).

% ------------------------------------------------------------
\section{Command-Line Interface}
\label{sec:cli}

\subsection{General form}

\begin{lstlisting}[style=shell]
dnipro analyze <netlist> [options]
\end{lstlisting}

\code{<netlist>} is the path to a SPICE-like netlist file (UTF-8 encoding).

\subsection{Options}

\begin{longtable}{p{0.32\textwidth}p{0.60\textwidth}}
\toprule
Option & Description \\
\midrule
\code{--domain \{time|laplace\}} &
    Select the symbolic domain for equation construction.
    Default: \code{time}. \\[4pt]
\code{--numeric} &
    Apply numeric substitution to all free parameters.
    Must be combined with \code{--param} for each symbolic parameter present in the circuit.
    Raises an error if any parameter is missing. \\[4pt]
\code{--param NAME=VALUE} &
    Supply or override a parameter value.
    \code{NAME} must be a valid identifier; \code{VALUE} follows the same grammar
    as netlist values (plain number, SPICE suffix, or single symbol).
    May be repeated. \\[4pt]
\code{--freq FREQ\_HZ} &
    Substitute \(s = j2\pi f\) in Laplace-domain mode.
    \code{FREQ\_HZ} is parsed as a netlist value. \\[4pt]
\code{--omega OMEGA} &
    Substitute \(s = j\omega\) in Laplace-domain mode.
    \code{OMEGA} is parsed as a netlist value. \\[4pt]
\code{--output \{text|json|latex\}} &
    Select the output format. Default: \code{text}. \\[4pt]
\code{--linear-system} &
    Export a derived classical linear MNA system \(A x = b\) instead of the
    residual equations. Available in both \code{time} and \code{laplace}
    domains for systems linear in the unknown vector. \\[4pt]
\code{--descriptor-system} &
    Export the time-domain descriptor system \(Gx + C\dot{x} = b\). Mutually
    exclusive with \code{--linear-system}; rejected in \code{laplace} domain. \\[4pt]
\code{--formulas} &
    Legacy alias; ignored in version~0.1.0. \\
\bottomrule
\end{longtable}

\subsection{Exit codes}

\begin{longtable}{p{0.15\textwidth}p{0.77\textwidth}}
\toprule
Code & Meaning \\
\midrule
\code{0} & Analysis completed successfully; output written to stdout. \\
\code{1} & An error occurred (parse error, missing parameter, file I/O error).
           The error message is written to stderr. \\
\bottomrule
\end{longtable}

\subsection{Extended examples}

\begin{lstlisting}[style=shell]
# Time-domain text output (default)
dnipro analyze circuit.cir

# Laplace domain, JSON output
dnipro analyze circuit.cir --domain laplace --output json

# Frequency substitution at 1 kHz
dnipro analyze filter.cir --domain laplace --freq 1k --output text

# Angular frequency substitution
dnipro analyze filter.cir --domain laplace --omega 6283.18 --output json

# Full numeric substitution
dnipro analyze rc.cir --domain time --numeric \
    --param Vin=5 --param R=1k --param C=100n

# Partial symbolic + LaTeX output
dnipro analyze amp.cir --domain laplace --output latex

# Classical linear MNA matrix/vector report
dnipro analyze amp.cir --domain laplace --linear-system --output latex

# Time-domain descriptor report
dnipro analyze rc.cir --descriptor-system --output text
\end{lstlisting}

% ------------------------------------------------------------
\section{Packaged Examples and Tutorials}
\label{sec:packaged-examples}

The repository includes small example circuits and report-generation scripts.
These files are not required by the library API, but they are useful for checking
the implemented conventions against concrete circuits.

\subsection{\texttt{pub1}: small-signal VCCS example}

The \code{pub1/} directory contains a compact small-signal equivalent circuit and
ready-made matrix reports:

\begin{longtable}{p{0.36\textwidth}p{0.56\textwidth}}
\toprule
File & Purpose \\
\midrule
\code{pub1/smallsig\_full.cir} &
    Full \ngspice{}-style source netlist with \code{.param}, \code{.control},
    and reporting commands. It is intended for \ngspice{}, not direct
    \project{} parsing. \\
\code{pub1/smallsig\_dnipro\_v1.cir} &
    \project{} v1-compatible version of the same circuit. Use this file with
    \code{dnipro analyze}. \\
\code{pub1/smallsig\_mna\_matrices.py} &
    Python helper that prints the residual system, \(A x = b\), descriptor
    form, JSON, and \LaTeX{} fragments. \\
\code{pub1/smallsig\_linear\_mna.tex} &
    \LaTeX{} fragment for the classical linear \(A x = b\) report. \\
\code{pub1/smallsig\_descriptor\_mna.tex} &
    \LaTeX{} fragment for \(Gx + C\dot{x} = b\). For this purely resistive
    small-signal circuit, \(C\) is the zero matrix. \\
\code{pub1/*\_document.tex}, \code{pub1/*\_document.pdf} &
    Complete wrapper documents and generated PDFs for the two \LaTeX{}
    fragments. \\
\bottomrule
\end{longtable}

Useful commands:

\begin{lstlisting}[style=shell]
dnipro analyze pub1/smallsig_dnipro_v1.cir --linear-system
dnipro analyze pub1/smallsig_dnipro_v1.cir --descriptor-system
python pub1/smallsig_mna_matrices.py
\end{lstlisting}

\subsection{MOSFET small-signal tutorial}

The \code{tutorials/} directory contains a MOSFET small-signal equivalent model
expressed as explicit linear small-signal elements. This is not a built-in
MOSFET device model; it is a circuit-level equivalent built from supported
\project{} v1 primitives.

\begin{longtable}{p{0.40\textwidth}p{0.52\textwidth}}
\toprule
File & Purpose \\
\midrule
\code{tutorials/mosfet\_small\_signal.cir} &
    Full \ngspice{} AC analysis netlist with \code{.param}, \code{.ac}, and
    \code{.control}. \\
\code{tutorials/mss.out} &
    Captured \ngspice{} output containing AC sweep tables for \code{v(D)} and
    \code{i(Vin)}. \\
\code{tutorials/mosfet\_small\_signal\_dnipro.py} &
    Builds the \project{} Laplace-domain residual system and classical
    \(A x = b\) matrix for the same explicit small-signal equivalent. Optional
    flags print \LaTeX{} or large symbolic transfer expressions. \\
\code{tutorials/mosfet\_small\_signal\_numeric\_dnipro.py} &
    Evaluates the extracted MNA matrix numerically at the same AC sweep
    frequencies as \code{mss.out} and compares \code{v(D)} and \code{i(Vin)}
    against \ngspice{}. \\
\bottomrule
\end{longtable}

Symbolic matrix construction:

\begin{lstlisting}[style=shell]
python tutorials/mosfet_small_signal_dnipro.py
python tutorials/mosfet_small_signal_dnipro.py --latex
python tutorials/mosfet_small_signal_dnipro.py --solve
\end{lstlisting}

Numerical comparison against \ngspice{} output:

\begin{lstlisting}[style=shell]
python tutorials/mosfet_small_signal_numeric_dnipro.py
python tutorials/mosfet_small_signal_numeric_dnipro.py --all
\end{lstlisting}

The numeric script substitutes the tutorial parameters, evaluates
\(s = j2\pi f\), solves \(A(f)x=b\) for every frequency point, and reports the
absolute difference between \project{} and \ngspice{} results. In the checked-in
\code{mss.out} comparison, the observed maximum absolute differences are on the
order of \(10^{-5}\) for \code{v(D)} and \(10^{-9}\) for \code{i(Vin)}, consistent
with printed \ngspice{} output precision.

% ------------------------------------------------------------
\section{Netlist Syntax}
\label{sec:syntax}

\subsection{General rules}

A netlist is a plain UTF-8 text file.
Each non-empty, non-comment line defines one circuit element or model.
Lines beginning with \code{*} are comments and are ignored.
Blank lines are ignored.
Fields on an element line are separated by whitespace.

\subsection{Node names}

The following names are ground aliases and are normalized to the
internal ground node:

\begin{center}
\code{0} \quad \code{GND} \quad \code{gnd}
\end{center}

All other node names must match the identifier pattern
\verb|[A-Za-z_][A-Za-z0-9_]*|.
Node names are case-sensitive (except the three ground aliases above).
The ground node is never included in the unknown vector.

\subsection{Component names}

A component name must match the pattern \verb|[A-Za-z][A-Za-z0-9_]*|.
The first character determines the element type (case-insensitive).
The original name is preserved in output.
Duplicate component names in the same netlist are rejected.

\subsection{Supported element lines}

\begin{longtable}{p{0.06\textwidth}p{0.37\textwidth}p{0.47\textwidth}}
\toprule
Type & Syntax & Meaning \\
\midrule
R & \code{Rname np nm value} & Linear resistor. \\
C & \code{Cname np nm value} & Linear capacitor. \\
L & \code{Lname np nm value} & Linear inductor. \\
V & \code{Vname np nm source} & Independent voltage source. \\
I & \code{Iname np nm source} & Independent current source. \\
E & \code{Ename np nm ncp ncm gain} & Voltage-controlled voltage source (VCVS). \\
G & \code{Gname np nm ncp ncm gm} & Voltage-controlled current source (VCCS). \\
F & \code{Fname np nm Vctrl gain} & Current-controlled current source (CCCS). \\
H & \code{Hname np nm Vctrl rm} & Current-controlled voltage source (CCVS). \\
D & \code{Dname anode cathode model} & Diode. \\
.model & \code{.model name D(Is=\ldots{} n=\ldots{} Vt=\ldots{})} & Diode model definition. \\
\bottomrule
\end{longtable}

\subsection{Values}

Version~1 accepts three kinds of values for component parameters:

\begin{enumerate}[noitemsep]
    \item \textbf{Plain numbers}: integers, decimals, or scientific notation.
          Examples: \code{1}, \code{2.5}, \code{1e-6}, \code{-3.3e3}.
    \item \textbf{Numbers with SPICE suffixes} (see Table~\ref{tab:suffixes}).
          Examples: \code{1k}, \code{100n}, \code{25.85m}, \code{4.7u}.
    \item \textbf{Single symbolic identifiers} matching
          \verb|[A-Za-z_][A-Za-z0-9_]*|.
          Examples: \code{R}, \code{C1val}, \code{Vin}, \code{gm}.
\end{enumerate}

Arbitrary expressions such as \code{R/2} or \code{C*(1+alpha)} are
intentionally rejected in version~1.
Python \code{eval()} is never used for parsing.
Numeric values are represented as exact \sympy{} \code{Rational} objects.

\begin{table}[H]
\centering
\caption{SPICE suffix scale factors}
\label{tab:suffixes}
\begin{tabular}{lll}
\toprule
Suffix & Scale & Example \\
\midrule
\code{f}   & \(10^{-15}\) & \code{1f}    \\
\code{p}   & \(10^{-12}\) & \code{1p}    \\
\code{n}   & \(10^{-9}\)  & \code{100n}  \\
\code{u}   & \(10^{-6}\)  & \code{1u}    \\
\code{m}   & \(10^{-3}\)  & \code{25.85m}\\
\code{k}   & \(10^{3}\)   & \code{1k}    \\
\code{meg} & \(10^{6}\)   & \code{1meg}  \\
\code{g}   & \(10^{9}\)   & \code{1g}    \\
\code{t}   & \(10^{12}\)  & \code{1t}    \\
\bottomrule
\end{tabular}
\end{table}

The suffix \code{m} means milli (\(10^{-3}\)).
Use \code{meg} for mega (\(10^{6}\)).
Uppercase \code{M} is not recognized in version~1.

\subsection{Diode model syntax}

\begin{lstlisting}[language=spice]
.model ModelName D(Is=value n=value Vt=value)
\end{lstlisting}

Exactly the three parameters \code{Is}, \code{n}, and \code{Vt} must be
present (space-separated, not comma-separated).
Each value follows the standard value grammar.
Duplicate model names in the same netlist are rejected.

% ------------------------------------------------------------
\section{Source Functions}
\label{sec:sources}

Independent voltage (\code{V}) and current (\code{I}) sources accept
a source specification as their last field(s).
All source functions are converted to \sympy{} expressions of time \(t\)
during parsing.

\subsection{DC (constant)}

\begin{lstlisting}[language=spice]
V1 in 0 value
V1 in 0 DC(value)
\end{lstlisting}

Generates the constant expression \(f(t) = \textit{value}\).

\subsection{SIN (sinusoidal)}

\begin{lstlisting}[language=spice]
V1 in 0 SIN(offset amplitude omega)
\end{lstlisting}

Generates:
\begin{equation}
    f(t) = \textit{offset} + \textit{amplitude}\cdot\sin(\omega\, t).
\end{equation}

All three arguments are required.

\subsection{PULSE}

\begin{lstlisting}[language=spice]
V1 in 0 PULSE(v1 v2 delay rise_time fall_time width period)
\end{lstlisting}

All seven arguments are required (space-separated).
Generates a periodic, trapezoid-shaped waveform using \sympy{} \code{Piecewise}
and \code{floor}.
Closed-form analytical solution of circuits driven by PULSE sources is not
required in version~1.

The periodic phase variable is:
\begin{equation}
    \tau = t - \textit{delay}
           - \textit{period}\left\lfloor
               \frac{t - \textit{delay}}{\textit{period}}
             \right\rfloor.
\end{equation}

\subsection{PWL (piecewise linear)}

\begin{lstlisting}[language=spice]
V1 in 0 PWL(t0 v0 t1 v1 t2 v2 ...)
\end{lstlisting}

At least two time/value pairs are required.
Arguments must be given in pairs (even count); values are space-separated.
Generates a \sympy{} \code{Piecewise} linear interpolation between the provided
points. The source holds \(v_0\) for \(t \leq t_0\) and \(v_{n-1}\) for
\(t > t_{n-1}\).

\subsection{SQUARE (square wave)}

\begin{lstlisting}[language=spice]
V1 in 0 SQUARE(offset amplitude omega duty)
\end{lstlisting}

All four arguments are required.
\code{duty} is the duty cycle fraction in \((0, 1)\).
Generates a periodic square wave using a normalized phase variable:
\begin{equation}
    \phi = \frac{\omega\,t}{2\pi}
           - \left\lfloor\frac{\omega\,t}{2\pi}\right\rfloor \in [0,1).
\end{equation}
\begin{equation}
    f(t) = \begin{cases}
        \textit{offset} + \textit{amplitude} & \phi < \textit{duty}, \\
        \textit{offset} - \textit{amplitude} & \text{otherwise.}
    \end{cases}
\end{equation}

\subsection{SAW (sawtooth)}

\begin{lstlisting}[language=spice]
V1 in 0 SAW(offset amplitude period)
\end{lstlisting}

All three arguments are required.
Generates a periodic linear ramp:
\begin{equation}
    \tau = t - \textit{period}\left\lfloor\frac{t}{\textit{period}}\right\rfloor,
    \qquad
    f(t) = \textit{offset} + \textit{amplitude}\,\frac{\tau}{\textit{period}}.
\end{equation}

% ============================================================
% PART II — MATHEMATICAL REFERENCE
% ============================================================
\part{Mathematical Reference}

% ------------------------------------------------------------
\section{MNA Theory}
\label{sec:mna-theory}

\subsection{Overview of Modified Nodal Analysis}

\mna{} is a systematic method for formulating circuit equations.
The standard MNA matrix form is
\begin{equation}
    G\,x(t) + C\,\dot{x}(t) = b(t),
\end{equation}
where \(G\) is the conductance matrix, \(C\) is the reactive (capacitive/inductive)
matrix, and \(b(t)\) collects source contributions.

\project{} does not use the matrix form as its primary representation.
Instead it constructs the equivalent \emph{residual form}:
\begin{equation}
    F\!\left(x(t),\,\dot{x}(t),\,t,\,p\right) = 0.
\end{equation}
The residual form is more general: it accommodates nonlinear components (such as
diodes) in the same framework without restructuring the solver.
Linear matrix forms are special cases of the residual form.

\subsection{Unknown vector}

For a circuit with \(N\) non-ground nodes and \(B\) voltage-defined branches
the unknown vector is:
\begin{equation}
    x(t) =
    \begin{bmatrix}
        v_{n_1}(t) \\ \vdots \\ v_{n_N}(t) \\
        i_{b_1}(t) \\ \vdots \\ i_{b_B}(t)
    \end{bmatrix}.
\end{equation}

In \emph{time-domain} mode, node voltages and branch currents are
\sympy{} \code{Function} objects applied to the time symbol \(t\),
e.g.\ \(v_{\mathrm{out}}(t)\) and \(i_{V1}(t)\).
This enables \sympy{} \code{diff()} to generate symbolic derivatives.

In \emph{Laplace-domain} mode, they are plain \sympy{} \code{Symbol} objects
e.g.\ \(v_{\mathrm{out}}\) and \(i_{V1}\), with the complex frequency
\(s = \texttt{Symbol("s")}\).

The ordering is deterministic:
\begin{enumerate}[noitemsep]
    \item Non-ground node voltages in first-appearance order.
    \item Branch current unknowns in first-appearance order of their components.
\end{enumerate}

\subsection{Residual vector}

The residual vector \(F\) has the same length as \(x\).
Its entries are ordered as:
\begin{enumerate}[noitemsep]
    \item One KCL residual per non-ground node, in node order.
    \item One branch constraint residual per voltage-defined branch, in
          component order.
\end{enumerate}

\subsection{KCL residuals}

For each non-ground node \(n\), the KCL residual is formed as the
algebraic sum of currents \emph{leaving} the node:
\begin{equation}
    F_n = \sum_k i_k^{(\mathrm{out})}(n) = 0.
\end{equation}
Here, \(i_k^{(\mathrm{out})}(n)\) is positive when current flows away from
node~\(n\) through branch~\(k\).

Internally, the builder maintains one accumulator per non-ground node.
The operation \code{add\_current(np,~nm,~current)} adds \code{current} to
node~\code{np} and subtracts it from node~\code{nm}.
If either terminal is the ground node, no residual is created for ground.

\subsection{Voltage-defined branches}

The following elements introduce a branch-current unknown and a branch
constraint residual:

\begin{center}
\begin{tabular}{ll}
\toprule
Element & Branch constraint \\
\midrule
Voltage source \code{V}  & \(v(np) - v(nm) - V_s = 0\) \\
VCVS \code{E}            & \(v(np) - v(nm) - \mu[v(ncp)-v(ncm)] = 0\) \\
CCVS \code{H}            & \(v(np) - v(nm) - r_m\,i_{ctrl} = 0\) \\
Inductor \code{L}        & \(v(np) - v(nm) - L\,\dot{i}_L = 0\) (time) \\
                         & \(v(np) - v(nm) - sL\,i_L = 0\) (Laplace) \\
\bottomrule
\end{tabular}
\end{center}

VCCS (\code{G}), CCCS (\code{F}), resistor (\code{R}), capacitor (\code{C}),
current source (\code{I}), and diode (\code{D}) do \emph{not} introduce
branch-current unknowns.

\subsection{Reference direction conventions}

For any two-terminal element written as \code{Xname np nm}:
\begin{equation}
    v_X = v(np) - v(nm),
    \qquad
    i_X: np \to nm \text{ (positive direction)}.
\end{equation}

For independent sources:
\begin{itemize}[noitemsep]
    \item A voltage source constrains \(v(np) - v(nm) = V_s\);
          its branch current flows from \code{np} to \code{nm}.
    \item A current source injects current \emph{from} \code{np} \emph{to} \code{nm}
          through the external circuit.
\end{itemize}

% ------------------------------------------------------------
\section{Component Equations}
\label{sec:component-equations}

\subsection{Resistor}

For \code{R1 np nm R}, the current from \code{np} to \code{nm} is:
\begin{equation}
    i_R(t) = \frac{v(np,t) - v(nm,t)}{R}.
\end{equation}
No branch-current unknown is allocated.
KCL contribution via \code{add\_current(np, nm, \(i_R\))}.

\subsection{Capacitor}

\textbf{Time domain.} For \code{C1 np nm C}:
\begin{equation}
    i_C(t) = C\,\frac{d}{dt}\bigl[v(np,t) - v(nm,t)\bigr].
\end{equation}
Implemented as \code{C * diff(voltage, t)}.

\textbf{Laplace domain.} The admittance representation is used:
\begin{equation}
    i_C = sC\,\bigl[v(np) - v(nm)\bigr].
\end{equation}
No branch-current unknown in either domain.

\subsection{Inductor}

An inductor is treated as a voltage-defined branch in both domains.

\textbf{Time domain.} For \code{L1 np nm L}:
Branch current \(i_{L1}(t)\) is allocated.
\begin{align}
    \text{KCL}: & \quad i_{L1}(t) \text{ added from \code{np} to \code{nm},} \\
    \text{Constraint}: & \quad v(np,t) - v(nm,t) - L\,\frac{di_{L1}(t)}{dt} = 0.
\end{align}

\textbf{Laplace domain.}
\begin{align}
    \text{KCL}: & \quad i_{L1} \text{ added from \code{np} to \code{nm},} \\
    \text{Constraint}: & \quad v(np) - v(nm) - sL\,i_{L1} = 0.
\end{align}

\subsection{Voltage source}

For \code{V1 np nm source}:
Branch current \(i_{V1}(t)\) is allocated.
\begin{align}
    \text{KCL}: & \quad i_{V1}(t) \text{ added from \code{np} to \code{nm},} \\
    \text{Constraint}: & \quad v(np,t) - v(nm,t) - f(t) = 0,
\end{align}
where \(f(t)\) is the source expression (DC, SIN, etc.).

\subsection{Current source}

For \code{I1 np nm source}:
The source expression \(f(t)\) is added directly to KCL as current from
\code{np} to \code{nm}.
No branch-current unknown is allocated.

\subsection{VCVS (E source)}

For \code{E1 np nm ncp ncm mu}:
Branch current \(i_{E1}(t)\) is allocated.
\begin{align}
    \text{KCL}: & \quad i_{E1}(t) \text{ added from \code{np} to \code{nm},} \\
    \text{Constraint}: & \quad
        v(np) - v(nm) - \mu\bigl[v(ncp) - v(ncm)\bigr] = 0.
\end{align}

\subsection{VCCS (G source)}

For \code{G1 np nm ncp ncm gm}:
The output current from \code{np} to \code{nm} is:
\begin{equation}
    i_G = g_m\,\bigl[v(ncp) - v(ncm)\bigr].
\end{equation}
No branch-current unknown is allocated.

\subsection{CCCS (F source)}

For \code{F1 np nm Vctrl beta}:
The output current from \code{np} to \code{nm} is:
\begin{equation}
    i_F = \beta\,i_{\mathrm{Vctrl}}.
\end{equation}
\code{Vctrl} must be the name of an existing voltage-defined branch.
No branch-current unknown is allocated for the \code{F} element itself.

\subsection{CCVS (H source)}

For \code{H1 np nm Vctrl rm}:
Branch current \(i_{H1}\) is allocated.
\begin{align}
    \text{KCL}: & \quad i_{H1} \text{ added from \code{np} to \code{nm},} \\
    \text{Constraint}: & \quad
        v(np) - v(nm) - r_m\,i_{\mathrm{Vctrl}} = 0.
\end{align}
\code{Vctrl} must be the name of an existing voltage-defined branch.

\subsection{Diode}

For \code{D1 anode cathode Model}:
The diode voltage is \(v_D = v(\text{anode}) - v(\text{cathode})\).
The diode current from anode to cathode is the Shockley equation:
\begin{equation}
    i_D = I_S\!\left(\exp\!\left(\frac{v_D}{n\,V_T}\right) - 1\right),
\end{equation}
where \(I_S\), \(n\), \(V_T\) come from the named \code{.model} definition.
The current \(i_D\) is added to KCL via \code{add\_current(anode, cathode, \(i_D\))}.
No branch-current unknown is allocated.

% ------------------------------------------------------------
\section{Time Domain and Laplace Domain}
\label{sec:domains}

\subsection{Time-domain mode}

Selected with \code{--domain time} (the default).

Node voltages and branch currents are \sympy{} \code{Function} objects evaluated at
\(t = \texttt{Symbol("t",~real=True)}\).
For example, a node \code{out} gives the unknown \(v_{\mathrm{out}}(t)\).
Time derivatives are formed with \texttt{sympy.diff(expr, t)}.

Capacitor currents take the form \(C\,\dot{v}(t)\);
inductor constraints include \(L\,\dot{i}_L(t)\).

\subsection{Laplace-domain mode}

Selected with \code{--domain laplace}.

Node voltages and branch currents are plain \sympy{} \code{Symbol} objects.
The complex frequency is \(s = \texttt{Symbol("s")}\).
Capacitors contribute \(sC\,V\) and inductors use the constraint
\(V - sL\,I = 0\).

\subsection{Frequency substitution}

Frequency substitution replaces \(s\) with a concrete complex value.
It is only available in Laplace-domain mode.

\begin{center}
\begin{tabular}{ll}
\toprule
Option & Substitution \\
\midrule
\code{--freq F}   & \(s \leftarrow j\,2\pi f\) \\
\code{--omega W}  & \(s \leftarrow j\,\omega\) \\
\bottomrule
\end{tabular}
\end{center}

After substitution, residuals are complex-valued \sympy{} expressions.
The substitution is recorded in \code{ResidualSystem.metadata["frequency\_substitution"]}.

% ------------------------------------------------------------
\section{Nonlinear Component: Diode}
\label{sec:diode}

\subsection{Design restriction}

Version~1 supports exactly one nonlinear component: the junction diode.
The diode is handled inside the same residual system as all linear elements.
There is no separate nonlinear backend.

\subsection{What is not implemented}

The following are explicitly out of scope for version~1:
\begin{itemize}[noitemsep]
    \item Jacobian matrix construction or export.
    \item Newton-Raphson or any other nonlinear solver.
    \item Automatic small-signal linearization.
    \item BJT or MOSFET models.
\end{itemize}

If \code{--numeric} is applied to a circuit containing a diode,
\project{} substitutes numerical parameter values into the symbolic residuals
but does \emph{not} attempt to solve the resulting nonlinear system.

\subsection{Model parameters}

The diode model requires exactly three parameters:

\begin{center}
\begin{tabular}{llp{0.5\textwidth}}
\toprule
Parameter & Symbol & Meaning \\
\midrule
\code{Is} & \(I_S\) & Saturation current (A) \\
\code{n}  & \(n\)   & Emission coefficient (dimensionless) \\
\code{Vt} & \(V_T\) & Thermal voltage (V); typically \(\approx 25.85\,\text{mV}\) \\
\bottomrule
\end{tabular}
\end{center}

\subsection{Example}

\begin{lstlisting}[language=spice]
V1 in 0 Vin
R1 in out R
D1 out 0 Dmod
.model Dmod D(Is=1e-14 n=1 Vt=25.85m)
\end{lstlisting}

The residual system (time domain) will contain:
\begin{align}
    F_{\mathrm{in}} &=
        \frac{v_{\mathrm{in}}(t) - v_{\mathrm{out}}(t)}{R} + i_{V1}(t) = 0, \\
    F_{\mathrm{out}} &=
        \frac{v_{\mathrm{out}}(t) - v_{\mathrm{in}}(t)}{R}
        + I_S\!\left(\exp\!\left(\frac{v_{\mathrm{out}}(t)}{n\,V_T}\right)-1\right) = 0, \\
    F_{V1} &= v_{\mathrm{in}}(t) - V_{\mathrm{in}} = 0.
\end{align}

% ============================================================
% PART III — LIBRARY API REFERENCE
% ============================================================
\part{Library API Reference}

% ------------------------------------------------------------
\section{Package Overview}
\label{sec:package-overview}

\subsection{Public API}

The top-level package \code{dnipro} exports three convenience names:

\begin{lstlisting}[language=Python]
from dnipro import parse_netlist_file, parse_netlist_text, build_equations
import dnipro
print(dnipro.__version__)   # "0.1.0"
\end{lstlisting}

Additional public module-level APIs are documented in the sections below.
Symbols not documented in this manual should be treated as internal.

\subsection{Package structure}

\begin{lstlisting}
dnipro/
  __init__.py            # public exports
  errors.py              # exception hierarchy
  parser/
    netlist_parser.py    # parse_netlist_text, parse_netlist_file
    value_parser.py      # parse_value
    source_parser.py     # parse_source, source function classes
  circuit/
    circuit.py           # Circuit
    components.py        # Component subclasses
    models.py            # DiodeModel
    nodes.py             # normalize_node, GROUND
  mna/
    equation_builder.py  # build_equations, ResidualSystemBuilder
    linear_extraction.py # LinearSystem, DescriptorSystem, extraction functions
    residual_system.py   # ResidualSystem
  analysis/
    substitutions.py     # substitution functions
  export/
    text_exporter.py     # residual, linear, and descriptor text exporters
    json_exporter.py     # residual, linear, and descriptor JSON exporters
    latex_exporter.py    # residual, linear, and descriptor LaTeX exporters
  cli/
    main.py              # main()
\end{lstlisting}

% ------------------------------------------------------------
\section{Parser Module}
\label{sec:api-parser}

\subsection{\texttt{parse\_netlist\_text}}
\label{api:parse_netlist_text}

\begin{apisig}
parse\_netlist\_text(text: str) -> Circuit
\end{apisig}

\textbf{Module:} \code{dnipro.parser.netlist\_parser} (also exported from \code{dnipro})

\textbf{Parameters:}
\begin{itemize}[noitemsep]
    \item \code{text} --- the netlist as a string.
\end{itemize}

\textbf{Returns:} A \code{Circuit} object.

\textbf{Raises:}
\begin{itemize}[noitemsep]
    \item \code{ParseError} --- on malformed element lines, invalid node or component names,
          duplicate component names.
    \item \code{UnsupportedElementError} --- on unsupported element type letters or
          unsupported directives.
    \item \code{UnsupportedValueError} --- on disallowed value syntax (e.g.\ expressions,
          comma-separated source arguments).
    \item \code{ModelError} --- on malformed \code{.model} lines, unsupported model types,
          missing or extra model parameters, duplicate model names.
\end{itemize}

Comment lines (starting with \code{*}) and blank lines are silently ignored.

\subsection{\texttt{parse\_netlist\_file}}
\label{api:parse_netlist_file}

\begin{apisig}
parse\_netlist\_file(path: str) -> Circuit
\end{apisig}

\textbf{Module:} \code{dnipro.parser.netlist\_parser} (also exported from \code{dnipro})

\textbf{Parameters:}
\begin{itemize}[noitemsep]
    \item \code{path} --- filesystem path to the netlist file (UTF-8).
\end{itemize}

\textbf{Returns:} A \code{Circuit} object.

\textbf{Raises:} Same exceptions as \code{parse\_netlist\_text}; also
\code{OSError} on file I/O failure.

Thin wrapper around \code{parse\_netlist\_text}: reads the file and delegates.

\subsection{\texttt{parse\_value}}
\label{api:parse_value}

\begin{apisig}
parse\_value(raw: str) -> sympy.Expr
\end{apisig}

\textbf{Module:} \code{dnipro.parser.value\_parser}

\textbf{Parameters:}
\begin{itemize}[noitemsep]
    \item \code{raw} --- a value string; leading/trailing whitespace is stripped.
\end{itemize}

\textbf{Returns:} A \sympy{} expression: \code{Rational} for numeric values,
\code{Symbol} for identifier values.

\textbf{Raises:} \code{UnsupportedValueError} on empty input or disallowed syntax.

Numeric values are represented as exact \sympy{} \code{Rational} objects
(no floating-point rounding).

\subsection{\texttt{parse\_source}}
\label{api:parse_source}

\begin{apisig}
parse\_source(raw: str) -> SourceFunction
\end{apisig}

\textbf{Module:} \code{dnipro.parser.source\_parser}

\textbf{Parameters:}
\begin{itemize}[noitemsep]
    \item \code{raw} --- source specification string,
          e.g.\ \code{"5"}, \code{"SIN(0 Vm omega)"}, \code{"PULSE(0 1 0 1n 1n 500n 1u)"}.
\end{itemize}

\textbf{Returns:} A frozen dataclass implementing the \code{SourceFunction} protocol.
Call \code{.to\_sympy(t)} to obtain the \sympy{} expression.

\textbf{Raises:} \code{UnsupportedValueError} on unsupported function names,
wrong argument counts, or comma-separated arguments.

The returned object's \code{.kind} attribute is one of
\code{"DC"}, \code{"SIN"}, \code{"PULSE"}, \code{"PWL"}, \code{"SQUARE"}, \code{"SAW"}.

% ------------------------------------------------------------
\section{Data Model}
\label{sec:api-data-model}

\subsection{\texttt{Circuit}}
\label{api:circuit}

\begin{apisig}
@dataclass \\
class Circuit: \\
\quad components: list[Component] \\
\quad node\_order: list[str] \\
\quad component\_names: set[str] \\
\quad models: dict[str, DiodeModel]
\end{apisig}

\textbf{Module:} \code{dnipro.circuit.circuit}

Mutable container built incrementally by the parser.
The \code{node\_order} list preserves first-appearance order (excluding ground).

\textbf{Methods:}

\begin{apisig}
add\_component(component: Component) -> None
\end{apisig}
Appends \code{component} to \code{components}, records its name in
\code{component\_names}, and extends \code{node\_order} with any new nodes.
Raises \code{ParseError} on a duplicate component name.

\begin{apisig}
add\_model(model: DiodeModel) -> None
\end{apisig}
Adds \code{model} to \code{models} keyed by \code{model.name}.
Raises \code{ModelError} on a duplicate model name.

\subsection{Component classes}
\label{api:components}

All component classes are frozen dataclasses in \code{dnipro.circuit.components}.
The base class is:

\begin{apisig}
@dataclass(frozen=True) \\
class Component: \\
\quad name: str   \# component name (preserved from netlist) \\
\quad np: str     \# positive node (normalized) \\
\quad nm: str     \# negative node (normalized)
\end{apisig}

Subclasses and their additional fields:

\begin{longtable}{p{0.18\textwidth}p{0.72\textwidth}}
\toprule
Class & Additional fields \\
\midrule
\code{Resistor}      & \code{value: sympy.Expr} \\
\code{Capacitor}     & \code{value: sympy.Expr} \\
\code{Inductor}      & \code{value: sympy.Expr} \\
\code{VoltageSource} & \code{source\_value: sympy.Expr}
                       (already converted to SymPy via \code{parse\_source}) \\
\code{CurrentSource} & \code{source\_value: sympy.Expr} \\
\code{Diode}         & \code{model\_name: str} \\
\code{VCVS}          & \code{ncp: str}, \code{ncm: str}, \code{gain: sympy.Expr} \\
\code{VCCS}          & \code{ncp: str}, \code{ncm: str},
                       \code{transconductance: sympy.Expr} \\
\code{CCCS}          & \code{control\_branch: str}, \code{gain: sympy.Expr} \\
\code{CCVS}          & \code{control\_branch: str},
                       \code{transresistance: sympy.Expr} \\
\bottomrule
\end{longtable}

The fields \code{ncp} and \code{ncm} of \code{VCVS} and \code{VCCS} are the
positive and negative control nodes (normalized).
The field \code{control\_branch} of \code{CCCS} and \code{CCVS} is the name
of the voltage-defined branch whose current is sensed.

\subsection{\texttt{DiodeModel}}
\label{api:diodemodel}

\begin{apisig}
@dataclass(frozen=True) \\
class DiodeModel: \\
\quad name: str \\
\quad saturation\_current: sympy.Expr \\
\quad emission\_coefficient: sympy.Expr \\
\quad thermal\_voltage: sympy.Expr
\end{apisig}

\textbf{Module:} \code{dnipro.circuit.models}

\subsection{\texttt{ResidualSystem}}
\label{api:residualsystem}

\begin{apisig}
@dataclass(frozen=True) \\
class ResidualSystem: \\
\quad domain: str \\
\quad unknowns: list[sympy.Expr] \\
\quad residuals: list[sympy.Expr] \\
\quad node\_order: list[str] \\
\quad branch\_current\_order: list[str] \\
\quad residual\_names: list[str] \\
\quad metadata: dict
\end{apisig}

\textbf{Module:} \code{dnipro.mna.residual\_system}

A frozen, immutable value object.
Analysis functions (\code{apply\_numeric\_substitutions},
\code{apply\_frequency\_substitution}) do not modify this object;
they return a new \code{ResidualSystem}.

\textbf{Fields:}
\begin{itemize}[noitemsep]
    \item \code{domain} --- \code{"time"} or \code{"laplace"}.
    \item \code{unknowns} --- ordered list of \sympy{} expressions,
          length \(N + B\) (nodes then branch currents).
    \item \code{residuals} --- ordered list of \sympy{} residual expressions,
          same length as \code{unknowns}.
    \item \code{node\_order} --- names of non-ground nodes in first-appearance order.
    \item \code{branch\_current\_order} --- names of voltage-defined branches
          in first-appearance order.
    \item \code{residual\_names} --- human-readable labels for each residual,
          e.g.\ \code{"KCL(out)"}, \code{"V1"}, \code{"L1"}.
    \item \code{metadata} --- dictionary recording analysis history.
          Keys used by the library: \code{"substitutions"} (from numeric substitution)
          and \code{"frequency\_substitution"} (from frequency substitution).
\end{itemize}

The \(i\)-th unknown corresponds to the \(i\)-th residual.

\subsection{Derived linear algebraic form}

For residual systems that are linear in the unknown vector, \project{} can derive
the classical matrix equation
\begin{equation}
    A x = b.
\end{equation}
The unknown vector \(x\) is exactly \code{ResidualSystem.unknowns}; no reordering
is performed during extraction. The row order of \(A\) and \(b\) is exactly the
order of \code{ResidualSystem.residuals}, and row labels are preserved through
\code{ResidualSystem.residual\_names}.

The extraction uses \sympy{} linear-equation utilities after replacing
time-domain unknown functions such as \(v_{\mathrm{out}}(t)\) with temporary
symbols. This makes the operation compatible with both time-domain and
Laplace-domain \project{} residual systems.

If a residual contains a nonlinear dependence on an unknown, for example the
exponential current of a diode, extraction raises \code{NonlinearSystemError}.

\subsection{Derived descriptor form}

For time-domain linear systems, \project{} can also derive
\begin{equation}
    Gx(t) + C\dot{x}(t) = b(t),
\end{equation}
where
\begin{equation}
    \dot{x}(t) =
    \begin{bmatrix}
        \frac{d}{dt}x_0(t) & \cdots & \frac{d}{dt}x_{n-1}(t)
    \end{bmatrix}^{T}.
\end{equation}
The derivative vector preserves the same order as \(x\). Capacitor terms
contribute to \(C\) through voltage derivatives, and time-domain inductor branch
constraints contribute through branch-current derivatives. Descriptor extraction
is intentionally rejected for Laplace-domain systems because Laplace residuals
are already algebraic in \(s\).

% ------------------------------------------------------------
\section{MNA Equation Builder}
\label{sec:api-builder}

\subsection{\texttt{build\_equations}}
\label{api:build_equations}

\begin{apisig}
build\_equations(circuit: Circuit, domain: str = "time") -> ResidualSystem
\end{apisig}

\textbf{Module:} \code{dnipro.mna.equation\_builder}
(also exported from \code{dnipro})

\textbf{Parameters:}
\begin{itemize}[noitemsep]
    \item \code{circuit} --- a \code{Circuit} object returned by the parser.
    \item \code{domain} --- \code{"time"} (default) or \code{"laplace"}.
\end{itemize}

\textbf{Returns:} A \code{ResidualSystem}.

\textbf{Raises:}
\begin{itemize}[noitemsep]
    \item \code{UnsupportedElementError} --- on an unsupported domain string
          or an unrecognized component type.
    \item \code{ControlBranchError} --- when a CCCS or CCVS references a
          control branch that is not a voltage-defined branch (i.e.\ the named
          component does not appear in \code{branch\_current\_order}).
    \item \code{ModelError} --- when a \code{Diode} component references a
          model name not present in \code{circuit.models}.
\end{itemize}

This is the primary entry point for equation construction.
Internally it creates a \code{ResidualSystemBuilder} and calls its \code{build()} method.

\subsection{\texttt{ResidualSystemBuilder}}
\label{api:builder}

\begin{apisig}
@dataclass \\
class ResidualSystemBuilder: \\
\quad circuit: Circuit \\
\quad domain: str \\
\quad time: sympy.Symbol        \# Symbol("t", real=True) \\
\quad frequency: sympy.Symbol   \# Symbol("s")
\end{apisig}

\textbf{Module:} \code{dnipro.mna.equation\_builder}

Lower-level class used internally by \code{build\_equations}.
Available for advanced use when components need to be added programmatically.

\textbf{Methods:}

\begin{apisig}
node\_voltage(node\_name: str) -> sympy.Expr
\end{apisig}
Returns \code{Integer(0)} for the ground node.
Returns a \code{Function} of \code{self.time} in time-domain mode,
or a \code{Symbol} in Laplace-domain mode.

\begin{apisig}
branch\_current(component\_name: str) -> sympy.Expr
\end{apisig}
Returns or creates the branch-current unknown for the named component.
Creates a \code{Function} of \code{self.time} or a \code{Symbol} depending
on \code{self.domain}.

\begin{apisig}
add\_current(np: str, nm: str, current\_expr: sympy.Expr) -> None
\end{apisig}
Adds \code{current\_expr} to the KCL accumulator of node \code{np}
and subtracts it from node \code{nm}.
Ground nodes are silently skipped.

\begin{apisig}
add\_constraint(name: str, expr: sympy.Expr) -> None
\end{apisig}
Appends a branch constraint residual \code{expr = 0} labelled \code{name}.

\begin{apisig}
build() -> ResidualSystem
\end{apisig}
Processes all components in the circuit and returns the completed
\code{ResidualSystem}.

% ------------------------------------------------------------
\section{Linear Extraction Module}
\label{sec:api-linear-extraction}

All functions are in \code{dnipro.mna.linear\_extraction}. They operate on an
already-built \code{ResidualSystem} and return derived matrix views; the residual
system remains the primary representation.

\subsection{\texttt{LinearSystem}}
\label{api:linearsystem}

\begin{apisig}
@dataclass(frozen=True) \\
class LinearSystem: \\
\quad matrix: sympy.Matrix \\
\quad unknown\_vector: sympy.Matrix \\
\quad rhs: sympy.Matrix \\
\quad residual\_names: list[str] \\
\quad metadata: dict
\end{apisig}

\textbf{Module:} \code{dnipro.mna.linear\_extraction}

Represents the derived algebraic system \(A x = b\).
\code{matrix} is \(A\), \code{unknown\_vector} is \(x\), and \code{rhs} is \(b\).
The metadata dictionary is copied from the source residual system and includes
\code{"domain"}.

\subsection{\texttt{DescriptorSystem}}
\label{api:descriptorsystem}

\begin{apisig}
@dataclass(frozen=True) \\
class DescriptorSystem: \\
\quad conductance\_matrix: sympy.Matrix \\
\quad dynamic\_matrix: sympy.Matrix \\
\quad unknown\_vector: sympy.Matrix \\
\quad derivative\_vector: sympy.Matrix \\
\quad rhs: sympy.Matrix \\
\quad residual\_names: list[str] \\
\quad metadata: dict
\end{apisig}

\textbf{Module:} \code{dnipro.mna.linear\_extraction}

Represents the derived descriptor system \(Gx + C\dot{x} = b\).
\code{conductance\_matrix} is \(G\), \code{dynamic\_matrix} is \(C\), and
\code{derivative\_vector} is \(\dot{x}\). Convenience aliases
\code{DescriptorSystem.G} and \code{DescriptorSystem.C} return the same matrices.

\subsection{\texttt{extract\_linear\_system}}
\label{api:extract_linear_system}

\begin{apisig}
extract\_linear\_system(system: ResidualSystem) -> LinearSystem
\end{apisig}

\textbf{Parameters:}
\begin{itemize}[noitemsep]
    \item \code{system} --- a residual system from \code{build\_equations} or an
          analysis function.
\end{itemize}

\textbf{Returns:} A \code{LinearSystem} preserving the source unknown order and
residual row labels.

\textbf{Raises:} \code{NonlinearSystemError} if any residual is nonlinear in the
unknown vector.

\subsection{\texttt{extract\_descriptor\_system}}
\label{api:extract_descriptor_system}

\begin{apisig}
extract\_descriptor\_system(system: ResidualSystem) -> DescriptorSystem
\end{apisig}

\textbf{Parameters:}
\begin{itemize}[noitemsep]
    \item \code{system} --- a time-domain residual system.
\end{itemize}

\textbf{Returns:} A \code{DescriptorSystem} preserving source unknown order,
derivative order, and residual row labels.

\textbf{Raises:}
\begin{itemize}[noitemsep]
    \item \code{NonlinearSystemError} if the system is not time-domain.
    \item \code{NonlinearSystemError} if any residual is nonlinear in \(x\) or
          \(\dot{x}\).
\end{itemize}

% ------------------------------------------------------------
\section{Analysis Module}
\label{sec:api-analysis}

All functions are in \code{dnipro.analysis.substitutions}.
None of them modify their input; each returns a new \code{ResidualSystem}.

\subsection{\texttt{apply\_numeric\_substitutions}}
\label{api:apply_numeric_substitutions}

\begin{apisig}
apply\_numeric\_substitutions( \\
\quad system: ResidualSystem, \\
\quad substitutions: Mapping[sympy.Symbol, sympy.Expr], \\
) -> ResidualSystem
\end{apisig}

\textbf{Parameters:}
\begin{itemize}[noitemsep]
    \item \code{system} --- the input residual system.
    \item \code{substitutions} --- a mapping from \sympy{} \code{Symbol} objects
          to replacement \sympy{} expressions.
          Typically built from CLI \code{--param} arguments via
          \code{parse\_param\_assignments}.
\end{itemize}

\textbf{Returns:} A new \code{ResidualSystem} with residuals simplified after
substitution. The \code{metadata["substitutions"]} key records the applied mapping.

\textbf{Raises:} \code{MissingParameterError} if any free symbol (other than
\code{t} and \code{s}) in the system is absent from \code{substitutions}.

Uses \code{sympy.simplify} after substitution.
The unknowns list is not modified.

\subsection{\texttt{apply\_frequency\_substitution}}
\label{api:apply_frequency_substitution}

\begin{apisig}
apply\_frequency\_substitution( \\
\quad system: ResidualSystem, \\
\quad *, \\
\quad freq\_hz: sympy.Expr | None = None, \\
\quad omega: sympy.Expr | None = None, \\
) -> ResidualSystem
\end{apisig}

\textbf{Parameters:}
\begin{itemize}[noitemsep]
    \item \code{system} --- must have \code{domain == "laplace"}.
    \item \code{freq\_hz} --- frequency in Hz; substitutes
          \(s \leftarrow j\,2\pi f\).
    \item \code{omega} --- angular frequency; substitutes
          \(s \leftarrow j\,\omega\).
\end{itemize}

Exactly one of \code{freq\_hz} or \code{omega} must be provided.

\textbf{Returns:} A new \code{ResidualSystem} with \(s\) replaced.
Records \code{metadata["frequency\_substitution"]} with keys
\code{"mode"}, \code{"value"}, and \code{"s"} (the substituted value).

\textbf{Raises:} \code{UnsupportedValueError} if the domain is not
\code{"laplace"} or if both or neither of \code{freq\_hz}/\code{omega} are given.

\subsection{\texttt{collect\_parameter\_symbols}}
\label{api:collect_parameter_symbols}

\begin{apisig}
collect\_parameter\_symbols(system: ResidualSystem) -> set[sympy.Symbol]
\end{apisig}

\textbf{Returns:} All free \sympy{} symbols found in \code{system.residuals} and
\code{system.unknowns}, excluding the reserved independent variables
\code{t} (\code{Symbol("t",~real=True)}) and \code{s} (\code{Symbol("s")}).

Useful for determining which \code{--param} values are required before
calling \code{apply\_numeric\_substitutions}.

\subsection{\texttt{parse\_param\_assignment}}
\label{api:parse_param_assignment}

\begin{apisig}
parse\_param\_assignment(raw: str) -> tuple[sympy.Symbol, sympy.Expr]
\end{apisig}

\textbf{Parameters:}
\begin{itemize}[noitemsep]
    \item \code{raw} --- a string of the form \code{"NAME=VALUE"}.
\end{itemize}

\textbf{Returns:} A tuple \code{(Symbol(NAME), parse\_value(VALUE))}.

\textbf{Raises:} \code{UnsupportedValueError} on missing \code{=},
invalid name, empty value.

\subsection{\texttt{parse\_param\_assignments}}
\label{api:parse_param_assignments}

\begin{apisig}
parse\_param\_assignments( \\
\quad raw\_assignments: Iterable[str], \\
) -> dict[sympy.Symbol, sympy.Expr]
\end{apisig}

\textbf{Parameters:}
\begin{itemize}[noitemsep]
    \item \code{raw\_assignments} --- iterable of \code{"NAME=VALUE"} strings
          (e.g.\ the list of \code{--param} arguments from the CLI).
\end{itemize}

\textbf{Returns:} A dictionary mapping each \code{Symbol} to its replacement value.

\textbf{Raises:} \code{UnsupportedValueError} on malformed assignments or duplicates.

% ------------------------------------------------------------
\section{Export Functions}
\label{sec:api-export}

\subsection{\texttt{export\_text}}
\label{api:export_text}

\begin{apisig}
export\_text(system: ResidualSystem) -> str
\end{apisig}

\textbf{Module:} \code{dnipro.export.text\_exporter}

Returns a human-readable multi-line string.
The format is:

\begin{lstlisting}
dnipro residual system
domain: <domain>
nodes: <n1>, <n2>, ...
branch currents: <c1>, <c2>, ...
unknowns:
  [0] <expr>
  ...
residuals:
  <name>: <expr> = 0
  ...
metadata:
  <key>: <value>
  ...
\end{lstlisting}

Expressions are serialized with \code{sympy.sstr()} for stable output.
The \code{metadata} block is omitted when the metadata dictionary is empty.

\subsection{\texttt{export\_json}}
\label{api:export_json}

\begin{apisig}
export\_json(system: ResidualSystem) -> str
\end{apisig}

\textbf{Module:} \code{dnipro.export.json\_exporter}

Returns a JSON string with 2-space indentation.
All \sympy{} expressions are serialized as strings via \code{sympy.sstr()}.

Top-level keys (in output order):

\begin{longtable}{p{0.30\textwidth}p{0.62\textwidth}}
\toprule
Key & Value \\
\midrule
\code{project}              & \code{"dnipro"} \\
\code{domain}               & \code{"time"} or \code{"laplace"} \\
\code{unknowns}             & list of expression strings \\
\code{residual\_names}      & list of label strings \\
\code{residuals}            & list of expression strings \\
\code{node\_order}          & list of node name strings \\
\code{branch\_current\_order} & list of branch name strings \\
\code{metadata}             & nested object; SymPy values serialized to strings \\
\bottomrule
\end{longtable}

\subsection{\texttt{export\_latex}}
\label{api:export_latex}

\begin{apisig}
export\_latex(system: ResidualSystem) -> str
\end{apisig}

\textbf{Module:} \code{dnipro.export.latex\_exporter}

Returns a \LaTeX{} source fragment (not a complete document).
Uses \code{sympy.latex()} for mathematical notation.

The fragment contains:
\begin{itemize}[noitemsep]
    \item \verb|\section*{dnipro residual system}| with domain.
    \item \verb|\subsection*{Unknowns}| with an \code{align*} environment
          listing \(x_i = \ldots\).
    \item \verb|\subsection*{Residual Equations}| with an \code{align*} environment
          listing \(F_i(\textit{name}) = \ldots = 0\).
\end{itemize}

Include it in a document with \verb|\input{output.tex}| or paste inline.

\subsection{Linear and descriptor text exporters}
\label{api:linear_descriptor_text_exporters}

\begin{apisig}
export\_linear\_text(system: LinearSystem) -> str \\
export\_descriptor\_text(system: DescriptorSystem) -> str
\end{apisig}

\textbf{Module:} \code{dnipro.export.text\_exporter}

Return human-readable reports for \(A x = b\) and \(Gx + C\dot{x} = b\).
Both reports include the equation form, row-name mapping, unknown vector, matrix
or matrices, and right-hand-side vector. Descriptor reports additionally include
the derivative vector.

\subsection{Linear and descriptor JSON exporters}
\label{api:linear_descriptor_json_exporters}

\begin{apisig}
export\_linear\_json(system: LinearSystem) -> str \\
export\_descriptor\_json(system: DescriptorSystem) -> str
\end{apisig}

\textbf{Module:} \code{dnipro.export.json\_exporter}

Return JSON strings with stable \sympy{} \code{sstr()} serialization. Linear
reports use keys \code{"form"}, \code{"unknowns"}, \code{"residual\_names"},
\code{"matrix"}, \code{"rhs"}, and \code{"metadata"}. Descriptor reports use
\code{"conductance\_matrix"}, \code{"dynamic\_matrix"}, and \code{"derivatives"}
in addition to the common keys.

\subsection{Linear and descriptor LaTeX exporters}
\label{api:linear_descriptor_latex_exporters}

\begin{apisig}
export\_linear\_latex(system: LinearSystem) -> str \\
export\_descriptor\_latex(system: DescriptorSystem) -> str
\end{apisig}

\textbf{Module:} \code{dnipro.export.latex\_exporter}

Return \LaTeX{} source fragments using \code{sympy.latex()} for matrix and vector
entries. Linear output contains row and unknown mapping tables plus displayed
\(A\), \(x\), and \(b\). Descriptor output contains row mapping and
unknown/derivative mapping tables plus displayed \(G\), \(C\), \(x\),
\(\dot{x}\), and \(b\).

% ------------------------------------------------------------
\section{Error Hierarchy}
\label{sec:errors}

All \project{}-specific exceptions inherit from \code{DniproError}.
Catching \code{DniproError} covers all library errors.

\begin{lstlisting}[language=Python]
DniproError(Exception)
  ParseError(DniproError)
    UnsupportedElementError(ParseError)
    UnsupportedValueError(ParseError)
  ModelError(DniproError)
  ControlBranchError(DniproError)
  MissingParameterError(DniproError)
  NonlinearSystemError(DniproError)
\end{lstlisting}

\begin{longtable}{p{0.33\textwidth}p{0.59\textwidth}}
\toprule
Exception & When raised \\
\midrule
\code{DniproError}           & Base class; never raised directly. \\
\code{ParseError}            & Malformed netlist line, invalid name, duplicate component. \\
\code{UnsupportedElementError} & Unknown element type letter or unsupported directive. \\
\code{UnsupportedValueError} & Disallowed value syntax, unsupported source function,
                               comma-separated arguments, wrong argument count. \\
\code{ModelError}            & Malformed \code{.model} line, unsupported model type,
                               missing/extra model parameters, duplicate model name. \\
\code{ControlBranchError}    & CCCS or CCVS control branch is not a voltage-defined branch. \\
\code{MissingParameterError} & \code{--numeric} requested but one or more parameters
                               are absent.
                               \code{exc.missing} is a \code{tuple[str,~\ldots{}]} of
                               missing symbol names. \\
\code{NonlinearSystemError}  & Linear or descriptor extraction requested for a system
                               that is nonlinear in the selected unknowns, or descriptor
                               extraction requested outside time-domain mode. \\
\bottomrule
\end{longtable}

% ============================================================
% PART IV — ARCHITECTURE AND DEVELOPER GUIDE
% ============================================================
\part{Architecture and Developer Guide}

% ------------------------------------------------------------
\section{Architecture Overview}
\label{sec:architecture}

\subsection{Architectural principles}

\begin{itemize}[noitemsep]
    \item The equation builder is the core of the system; the CLI is a thin wrapper.
    \item Parsers produce typed, immutable dataclass objects, not raw strings.
    \item Component contribution is implemented via direct method calls on
          the builder, not through virtual matrix stamps.
    \item The primary internal representation is the residual system,
          not a coefficient matrix.
    \item No feature uses Python \code{eval()} on user-provided input.
    \item \ngspice{} is used only in external verification tests,
          never at runtime.
\end{itemize}

\subsection{Data flow}

\begin{figure}[H]
\centering
\begin{tikzpicture}[
    node distance=11mm,
    box/.style={draw, rounded corners, align=center,
                minimum width=42mm, minimum height=9mm, font=\small},
    arrow/.style={-{Stealth[length=2mm]}, thick}
]
\node[box] (netlist)  {SPICE-like netlist\\\texttt{(.cir file)}};
\node[box] (parser)   [below=of netlist]  {Parser\\(\texttt{netlist\_parser})};
\node[box] (circuit)  [below=of parser]   {Circuit data model\\(\texttt{Circuit, Component})};
\node[box] (builder)  [below=of circuit]  {Equation builder\\(\texttt{ResidualSystemBuilder})};
\node[box] (system)   [below=of builder]  {Residual system\\$F(x,\dot{x},t,p)=0$};

\node[box] (analysis) [below left=10mm and 20mm of system]
                                           {Analysis\\(\texttt{substitutions})};
\node[box] (linear)   [below=10mm of system]
                                           {Linear extraction\\$Ax=b$, $Gx+C\dot{x}=b$};
\node[box] (export)   [below right=10mm and 20mm of system]
                                           {Export\\(text / JSON / \LaTeX{})};

\draw[arrow] (netlist)  -- (parser);
\draw[arrow] (parser)   -- (circuit);
\draw[arrow] (circuit)  -- (builder);
\draw[arrow] (builder)  -- (system);
\draw[arrow] (system)   -- (analysis);
\draw[arrow] (system)   -- (linear);
\draw[arrow] (system)   -- (export);
\draw[arrow] (linear)   -- (export);
\draw[arrow] (analysis) -- ++(0,-10mm)
    node[below, font=\small] {new \texttt{ResidualSystem}};
\end{tikzpicture}
\caption{dnipro data flow.}
\label{fig:dataflow}
\end{figure}

\subsection{Module dependency graph}

\begin{center}
\begin{tabular}{lp{0.6\textwidth}}
\toprule
Module & Imports from \\
\midrule
\code{cli.main}            & \code{parser}, \code{mna}, \code{analysis}, \code{export} \\
\code{mna.equation\_builder}  & \code{circuit.*}, \code{mna.residual\_system}, \code{errors} \\
\code{mna.linear\_extraction} & \code{mna.residual\_system}, \code{errors} \\
\code{parser.netlist\_parser} & \code{circuit.*}, \code{parser.value\_parser},
                               \code{parser.source\_parser}, \code{errors} \\
\code{analysis.substitutions} & \code{mna.residual\_system}, \code{parser.value\_parser},
                               \code{errors} \\
\code{export.*}            & \code{mna.residual\_system}, \code{mna.linear\_extraction} \\
\code{circuit.*}           & \code{errors} \\
\code{errors}              & (none) \\
\bottomrule
\end{tabular}
\end{center}

% ------------------------------------------------------------
\section{Adding a New Component}
\label{sec:extending}

To add a new two-terminal linear component \code{X} to version~1:

\begin{enumerate}

\item \textbf{Add a frozen dataclass} in \code{dnipro/circuit/components.py}:
\begin{lstlisting}[language=Python]
@dataclass(frozen=True)
class XElement(Component):
    value: sp.Expr
\end{lstlisting}

\item \textbf{Add a parser branch} in \code{dnipro/parser/netlist\_parser.py}:
Add an \code{if element\_type == "X":} branch in \code{\_parse\_element\_line}
and a helper \code{\_parse\_xelement} function following the pattern of
\code{\_parse\_resistor}.

\item \textbf{Implement the contribution method} in
\code{dnipro/mna/equation\_builder.py}:
Add \code{\_contribute\_xelement(self, component: XElement)} following the
equations for the component.
Use \code{self.add\_current(np, nm, expr)} for KCL contributions and
\code{self.add\_constraint(name, expr)} for branch constraints.

\item \textbf{Dispatch in \texttt{build()}}:
Add an \code{elif isinstance(component, XElement):} branch in the
\code{build()} method to call \code{\_contribute\_xelement}.

\item \textbf{Write tests} covering:
\begin{itemize}[noitemsep]
    \item Parser: valid line, field-count errors, invalid values.
    \item MNA builder: correct KCL residuals and constraints in both domains.
    \item Use \code{sympy.simplify(actual - expected) == 0} for equation checks.
\end{itemize}

\end{enumerate}

Voltage-defined branches additionally require:
pre-allocating the branch current in the first pass of \code{build()}
by adding \code{XElement} to the \code{isinstance} check:
\begin{lstlisting}[language=Python]
if isinstance(component, (VoltageSource, VCVS, CCVS, Inductor, XElement)):
    self.branch_current(component.name)
\end{lstlisting}

% ------------------------------------------------------------
\section{Testing Guide}
\label{sec:testing}

\subsection{Test structure}

\begin{lstlisting}
tests/
  parser/
    test_value_parser.py
    test_source_parser.py
    test_netlist_parser_rvi.py
    test_netlist_parser_cl.py
    test_netlist_parser_dependent_sources.py
    test_netlist_parser_diode.py
  mna/
    test_equation_builder_rvi.py
    test_equation_builder_cl.py
    test_equation_builder_dependent_sources.py
    test_equation_builder_diode.py
    test_equation_builder_time_sources.py
    test_equation_builder_laplace.py
    test_linear_extraction.py
    test_reference_circuit_coverage.py
  analysis/
    test_substitutions.py
  export/
    test_exporters.py
    test_linear_system_exporters.py
  circuit/
    test_nodes.py
  cli/
    test_main.py
  integration/
    test_ngspice_verification.py
\end{lstlisting}

\subsection{Mathematical equivalence}

Tests compare symbolic expressions using \sympy{} simplification:
\begin{lstlisting}[language=Python]
import sympy as sp

def assert_equivalent(actual: sp.Expr, expected: sp.Expr) -> None:
    assert sp.simplify(actual - expected) == 0
\end{lstlisting}

Output-format tests (JSON, text, LaTeX) may use exact string matching
where stable formatting is intentionally verified.

\subsection{ngspice integration tests}

Tests in \code{tests/integration/} are decorated with
\code{@pytest.mark.ngspice} and skipped when \ngspice{} is absent:
\begin{lstlisting}[language=Python]
import shutil, pytest

@pytest.mark.ngspice
@pytest.mark.skipif(
    shutil.which("ngspice") is None,
    reason="ngspice is not installed"
)
def test_dc_divider():
    ...
\end{lstlisting}

Run them explicitly with \code{pytest -m ngspice}.

\subsection{Current local health check}

At the time this manual was updated, the project health check used:

\begin{lstlisting}[style=shell]
/home/konst/Documents/Projects/venv/bin/python -m pytest -q
/home/konst/Documents/Projects/venv/bin/ruff check .
/home/konst/Documents/Projects/venv/bin/python -m pytest --collect-only -q
\end{lstlisting}

The collected test count is 165 tests. The full test suite and Ruff check pass
in the current repository state.

% ============================================================
% Appendix
% ============================================================
\appendix

\section{Summary of Public API}
\label{app:api-summary}

\begin{longtable}{p{0.45\textwidth}p{0.47\textwidth}}
\toprule
Signature & Module \\
\midrule
\code{parse\_netlist\_text(text: str) -> Circuit}
    & \code{dnipro} / \code{dnipro.parser.netlist\_parser} \\
\code{parse\_netlist\_file(path: str) -> Circuit}
    & \code{dnipro} / \code{dnipro.parser.netlist\_parser} \\
\code{build\_equations(circuit, domain="time") -> ResidualSystem}
    & \code{dnipro} / \code{dnipro.mna.equation\_builder} \\
\code{extract\_linear\_system(system) -> LinearSystem}
    & \code{dnipro.mna.linear\_extraction} \\
\code{extract\_descriptor\_system(system) -> DescriptorSystem}
    & \code{dnipro.mna.linear\_extraction} \\
\code{export\_text(system) -> str}
    & \code{dnipro.export.text\_exporter} \\
\code{export\_json(system) -> str}
    & \code{dnipro.export.json\_exporter} \\
\code{export\_latex(system) -> str}
    & \code{dnipro.export.latex\_exporter} \\
\code{export\_linear\_text(system) -> str}
    & \code{dnipro.export.text\_exporter} \\
\code{export\_linear\_json(system) -> str}
    & \code{dnipro.export.json\_exporter} \\
\code{export\_linear\_latex(system) -> str}
    & \code{dnipro.export.latex\_exporter} \\
\code{export\_descriptor\_text(system) -> str}
    & \code{dnipro.export.text\_exporter} \\
\code{export\_descriptor\_json(system) -> str}
    & \code{dnipro.export.json\_exporter} \\
\code{export\_descriptor\_latex(system) -> str}
    & \code{dnipro.export.latex\_exporter} \\
\code{apply\_numeric\_substitutions(system, subs) -> ResidualSystem}
    & \code{dnipro.analysis.substitutions} \\
\code{apply\_frequency\_substitution(system, *, freq\_hz, omega) -> ResidualSystem}
    & \code{dnipro.analysis.substitutions} \\
\code{collect\_parameter\_symbols(system) -> set[Symbol]}
    & \code{dnipro.analysis.substitutions} \\
\code{parse\_param\_assignment(raw) -> tuple[Symbol, Expr]}
    & \code{dnipro.analysis.substitutions} \\
\code{parse\_param\_assignments(raws) -> dict[Symbol, Expr]}
    & \code{dnipro.analysis.substitutions} \\
\code{parse\_value(raw: str) -> sympy.Expr}
    & \code{dnipro.parser.value\_parser} \\
\code{parse\_source(raw: str) -> SourceFunction}
    & \code{dnipro.parser.source\_parser} \\
\bottomrule
\end{longtable}

\section{Reference Circuit Index}
\label{app:reference-circuits}

The following reference circuits are defined in the technical specification
and are covered by tests in \code{tests/mna/test\_reference\_circuit\_coverage.py}
and \code{tests/integration/test\_ngspice\_verification.py}.

\begin{longtable}{p{0.12\textwidth}p{0.28\textwidth}p{0.50\textwidth}}
\toprule
ID & Circuit & Test purpose \\
\midrule
TDD-CIR-001 & Voltage divider          & Basic KCL, voltage source constraint. \\
TDD-CIR-002 & RC low-pass              & Capacitor time derivative. \\
TDD-CIR-003 & RC high-pass             & Floating capacitor, non-ground nodes. \\
TDD-CIR-004 & RL low-pass              & Inductor branch current. \\
TDD-CIR-005 & RLC low-pass             & Second-order dynamics. \\
TDD-CIR-006 & Current-driven RC        & Current source sign convention. \\
TDD-CIR-007 & VCVS amplifier           & E source constraint. \\
TDD-CIR-008 & VCCS transconductance    & G source current direction. \\
TDD-CIR-009 & CCCS current gain        & F source, sensing branch. \\
TDD-CIR-010 & CCVS transresistance     & H source constraint. \\
TDD-CIR-011 & SIN voltage source RC    & Time-dependent V source. \\
TDD-CIR-012 & SIN current source RC    & Time-dependent I source. \\
TDD-CIR-013 & Pulse-driven RC          & PULSE source formula. \\
TDD-CIR-014 & PWL-driven RC            & PWL parser and formula. \\
TDD-CIR-015 & Diode with resistor      & Nonlinear residual formula. \\
\bottomrule
\end{longtable}

\section{Version History}
\label{app:changelog}

\begin{longtable}{p{0.15\textwidth}p{0.75\textwidth}}
\toprule
Version & Changes \\
\midrule
0.1.0 & Initial release.
        Supported elements: R, C, L, V, I, E, G, F, H, D.
        Source functions: DC, SIN, PULSE, PWL, SQUARE, SAW.
        Domains: time, Laplace.
        Exports: text, JSON, LaTeX.
        Linear reporting: residual form, \(A x = b\), and descriptor
        \(Gx + C\dot{x} = b\).
        CLI: \texttt{dnipro analyze}, \texttt{--linear-system},
        \texttt{--descriptor-system}.
        Documentation includes \code{pub1} matrix reports and MOSFET
        small-signal tutorial verification against \ngspice{} output. \\
\bottomrule
\end{longtable}

\end{document}
